\documentclass[11pt]{article}
\usepackage[margin=1in]{geometry}
\usepackage{graphicx}
\usepackage{subcaption}
\usepackage{minted}
\usepackage{booktabs}
\usepackage{float}
\usepackage{hyperref}

\graphicspath{{figures/}}

% \onefig{file}{caption}
\newcommand{\onefig}[2]{%
  \begin{figure}[H]
    \centering
    \includegraphics[width=0.6\textwidth]{#1}
    \caption{#2}
\end{figure}}

% \twofig{left file}{left caption}{right file}{right caption}{overall caption}
\newcommand{\twofig}[5]{%
  \begin{figure}[H]
    \centering
    \begin{subfigure}{0.48\textwidth}
      \includegraphics[width=\linewidth]{#1}
      \caption{#2}
    \end{subfigure}\hfill
    \begin{subfigure}{0.48\textwidth}
      \includegraphics[width=\linewidth]{#3}
      \caption{#4}
    \end{subfigure}
    \caption{#5}
\end{figure}}

% \tableinput{file}{caption}
\newcommand{\tableinput}[2]{%
  \begin{table}[H]
    \centering
    \caption{#2}
    \input{tables/#1}
\end{table}}

\title{CS 4372 Assignment 3: Intel Image Classification with Transfer Learning}
\author{}
\date{\today}

\begin{document}
\maketitle

% =====================================================================
\section{Data}

The data set is the Intel Image Classification data set from Kaggle,
which contains natural scene images of size $150 \times 150$ pixels
split into six classes: buildings, forest, glacier, mountain, sea, and
street. The data comes pre-split into a training set
(\texttt{seg\_train}, 14{,}034 images), a labeled test set
(\texttt{seg\_test}, 3{,}000 images), and an unlabeled prediction set
(\texttt{seg\_pred}, 7{,}301 images). Since the prediction set has no
labels, the labeled test set is used as the validation set throughout
this report. Table~\ref{tab:counts} shows that the classes are roughly
balanced, so accuracy is a reasonable summary metric.

Images were loaded with Keras' \texttt{image\_dataset\_from\_directory}
at their native $150 \times 150$ resolution, shuffled, and batched in
groups of 32. Pixel values are passed through the ResNet50
\texttt{preprocess\_input} function inside the model, which converts
RGB to BGR and subtracts the ImageNet channel means.

\begin{table}[H]
  \centering
  \caption{Number of images per class.}
  \label{tab:counts}
  \begin{tabular}{lrr}
\toprule
Class & Train & Validation \\
\midrule
buildings & 2191 & 437 \\
forest & 2271 & 474 \\
glacier & 2404 & 553 \\
mountain & 2512 & 525 \\
sea & 2274 & 510 \\
street & 2382 & 501 \\
\midrule
Total & 14034 & 3000 \\
\bottomrule
\end{tabular}

\end{table}
\onefig{sample_images}{A batch of training images with their labels.}

% =====================================================================
\section{Base Model}

The model uses transfer learning from ResNet50 pre-trained on
ImageNet, with its classification head removed
(\texttt{include\_top=False}). The convolutional base produces a
$5 \times 5 \times 2048$ feature map, which is reduced to a
2048-dimensional vector by global average pooling, followed by a
dropout layer with rate 0.2 and a dense softmax layer with six
outputs. In this first stage the entire ResNet50 base is frozen and
run in inference mode (\texttt{training=False}), so its batch
normalization statistics are kept fixed and only the final dense layer
is trained.

The model was trained for 10 epochs with the Adam optimizer at a
learning rate of $10^{-4}$, minimizing sparse categorical
cross-entropy.

\tableinput{architecture}{Model architecture.}
\tableinput{base_model_results}{Base model accuracy and loss before
training and after 10 epochs.}
\twofig{base_accuracy}{Accuracy}{base_loss}{Cross-entropy loss}%
{Base model (frozen ResNet50): training and validation curves.}

% TODO: analysis of base model results

% =====================================================================
\section{Fine Tuning}

After training the classification head, the upper layers of the
ResNet50 base were unfrozen and trained together with the head using
the RMSprop optimizer at a lower learning rate. Two hyperparameters
were searched over: the number of base layers kept frozen,
$\{100, 110, 120\}$ out of 175, and the learning rate, $\{10^{-5},
10^{-6}, 10^{-7}\}$ (the base learning rate divided by 10, 100, and
1000). Each of the 9 combinations was trained for a single epoch.
Note that the combinations were run one after another on the same
model, so each starts from the weights left by the previous one rather
than from the base model.

The combination with the highest validation accuracy after its epoch
was selected: 120 frozen layers with a learning rate of $10^{-7}$
(Table~\ref{tab:chosen}).

The chosen combination was then trained for 10 more epochs (epochs
11--20).

\tableinput{grid_search}{Fine-tuning search: results after one epoch
for each combination.}
\tableinput{chosen_params}{Fine-tuning configuration used for the
final model.\label{tab:chosen}}
\onefig{fine_tune_curves}{Training and validation accuracy and loss
over the base model and fine-tuning stages.}

\subsection{Final Evaluation}
\tableinput{final_metrics}{Final model: per-class metrics on the
validation set.}
\onefig{confusion_matrix}{Final model: confusion matrix on the
validation set.}
\begin{figure}[H]
  \centering
  \includegraphics[width=0.9\textwidth]{prediction_grid}
  \caption{Final model: predicted and actual labels for 25 images from
    the labeled test set (\texttt{seg\_test}). Correct predictions are in
  green, incorrect in red.}
\end{figure}

\pagebreak
\section{Code}
\begin{minted}{python}
import itertools
import os

import matplotlib.pyplot as plt
import numpy as np
import tensorflow as tf
from numpy.random import shuffle

_URL = "https://www.kaggle.com/api/v1/datasets/download/puneet6060/intel-image-classification"
path_to_zip = tf.keras.utils.get_file("intel.zip", origin=_URL, extract=True)
PATH = os.path.join(os.path.dirname(path_to_zip), "intel_extracted")
train_dir = os.path.join(PATH, "seg_train/seg_train")
validation_dir = os.path.join(PATH, "seg_test/seg_test")
test_dir = os.path.join(PATH, "seg_pred/")

BATCH_SIZE = 32
IMG_SIZE = (150, 150)

train_dataset = tf.keras.utils.image_dataset_from_directory(
    train_dir, shuffle=True, batch_size=BATCH_SIZE, image_size=IMG_SIZE
)

validation_dataset = tf.keras.utils.image_dataset_from_directory(
    validation_dir, shuffle=True, batch_size=BATCH_SIZE, image_size=IMG_SIZE
)

test_dataset = tf.keras.utils.image_dataset_from_directory(
    test_dir, shuffle=True, batch_size=BATCH_SIZE, image_size=IMG_SIZE
)

preprocess_input = tf.keras.applications.resnet50.preprocess_input

AUTOTUNE = tf.data.AUTOTUNE
train_dataset = train_dataset.prefetch(buffer_size=AUTOTUNE)
validation_dataset = validation_dataset.prefetch(buffer_size=AUTOTUNE)
test_dataset = test_dataset.prefetch(buffer_size=AUTOTUNE)

IMG_SHAPE = IMG_SIZE + (3,)

base_model = tf.keras.applications.ResNet50(
    input_shape=IMG_SHAPE,
    include_top=False,
    weights="imagenet",
)

image_batch, label_batch = next(iter(train_dataset))
feature_batch = base_model(image_batch)

base_model.trainable = False
global_average_layer = tf.keras.layers.GlobalAveragePooling2D()
feature_batch_average = global_average_layer(feature_batch)

prediction_layer = tf.keras.layers.Dense(6, activation="softmax")
prediction_batch = prediction_layer(feature_batch_average)

inputs = tf.keras.Input(shape=IMG_SHAPE)
x = preprocess_input(inputs)
x = base_model(x, training=False)
x = global_average_layer(x)
x = tf.keras.layers.Dropout(0.2)(x)
outputs = prediction_layer(x)
model = tf.keras.Model(inputs, outputs)
model.summary()
base_learning_rate = 0.0001
model.compile(
    optimizer=tf.keras.optimizers.Adam(learning_rate=base_learning_rate),
    loss="sparse_categorical_crossentropy",
    metrics=["accuracy"],
)

initial_epochs = 10

loss0, accuracy0 = model.evaluate(validation_dataset)

history = model.fit(
    train_dataset, epochs=initial_epochs, validation_data=validation_dataset
)

history.history

acc = history.history["accuracy"]
val_acc = history.history["val_accuracy"]

loss = history.history["loss"]
val_loss = history.history["val_loss"]

plt.figure(figsize=(8, 8))
plt.plot(acc, label="Training Accuracy")
plt.plot(val_acc, label="Validation Accuracy")
plt.legend(loc="lower right")
plt.ylabel("Accuracy")
plt.ylim([min(plt.ylim()), 1])
plt.title("Training and Validation Accuracy")
plt.show()

plt.plot(loss, label="Training Loss")
plt.plot(val_loss, label="Validation Loss")
plt.legend(loc="upper right")
plt.ylabel("Cross Entropy")
plt.ylim([0, 10.0])
plt.title("Training and Validation Loss")
plt.xlabel("epoch")
plt.show()

## Fine Tuning
import itertools

base_model.trainable = True

len(base_model.layers)

fine_tune_at = [100 + (10 * x) for x in range(3)]
learning_rate_divider = [10**x for x in range(1, 4)]

combinations = itertools.product(fine_tune_at, learning_rate_divider)

curr_best = [None, 0]
histories = []

for f, r in combinations:
    for layer in base_model.layers[:f]:
        layer.trainable = False

    model.compile(
        optimizer=tf.keras.optimizers.RMSprop(learning_rate=base_learning_rate / r),
        loss="sparse_categorical_crossentropy",
        metrics=["accuracy"],
    )

    fine_tune_epochs = 1

    total_epochs = initial_epochs + fine_tune_epochs

    history_fine = model.fit(
        train_dataset,
        epochs=total_epochs,
        initial_epoch=len(history.epoch),
        validation_data=validation_dataset,
    )
    histories.append(history_fine)

    if curr_best[1] < history_fine.history["val_accuracy"][-1]:
        curr_best[0] = (f, r)
        curr_best[1] = history_fine.history["val_accuracy"][-1]

best_f, best_r = curr_best[0]

for layer in base_model.layers[:best_f]:
    layer.trainable = False

model.compile(
    optimizer=tf.keras.optimizers.RMSprop(learning_rate=base_learning_rate / best_r),
    loss="sparse_categorical_crossentropy",
    metrics=["accuracy"],
)

fine_tune_epochs = 10
total_epochs = initial_epochs + fine_tune_epochs

history_fine = model.fit(
    train_dataset,
    epochs=total_epochs,
    initial_epoch=len(history.epoch),
    validation_data=validation_dataset,
)

acc += history_fine.history["accuracy"]
val_acc += history_fine.history["val_accuracy"]

loss += history_fine.history["loss"]
val_loss += history_fine.history["val_loss"]

plt.figure(figsize=(8, 8))
plt.subplot(2, 1, 1)
plt.plot(acc, label="Training Accuracy")
plt.plot(val_acc, label="Validation Accuracy")
plt.ylim([0.8, 1])
plt.plot(
    [initial_epochs - 1, initial_epochs - 1], plt.ylim(), label="Start Fine Tuning"
)
plt.legend(loc="lower right")
plt.title("Training and Validation Accuracy")

plt.subplot(2, 1, 2)
plt.plot(loss, label="Training Loss")
plt.plot(val_loss, label="Validation Loss")
plt.ylim([0, 1.0])
plt.plot(
    [initial_epochs - 1, initial_epochs - 1], plt.ylim(), label="Start Fine Tuning"
)
plt.legend(loc="upper right")
plt.title("Training and Validation Loss")
plt.xlabel("epoch")
plt.show()

\end{minted}
\end{document}
